\section*{Bài 2:}
Khảo sát sự biến thiên và vẽ đồ thị hàm số: $y = -x^3 + 3x^2 - 4$

\subsection*{Lời giải:}
\begin{itemize}
\item[$\ast$] Tập xác định D=R
\item[$\ast$] Chiều biến thiên:
Ta có: $y^{\prime} = -3x^2 + 6x = -3x(x-2)$ \\
Xét phương trình $y^{\prime} = 0 \Leftrightarrow -3x(x-2) = 0 \Leftrightarrow x = 0$ hoặc $x = 2$
\item[$\ast$] Bảng biến thiên:
	\begin{figure}[H]
		\centering
		\begin{tikzpicture}
			\tkzTabInit[nocadre=false,lgt=2,espcl=2]
			{$x$ /1,$y'$ /1,$y$ /2}
			{$-\infty$,$0$,$2$,$+\infty$}
			\tkzTabLine{,-,$0$,+,$0$,-,}
			\tkzTabVar{+/ $+\infty$ ,-/$-4$,+/$0$,-/$-\infty$}
		\end{tikzpicture}
	\end{figure}
	Hàm số nghịch biến trên các khoảng $(-\infty; 0)$ và $(2; +\infty)$, đồng biến trên khoảng $(0; 2)$ \\
	Hàm số đạt cực đại tại điểm $x=2$, giá trị cực đại của hảm số là $y(2)=0$. \\
	Hàm số đạt cực tiểu tại điểm $=0$, giá trị cực tiểu của hàm số là $y(0)=-4$ \\
	Giới hạn của hàm số tại vô cực: 
	$$\lim_{x \to -\infty}{y} = +\infty$$
	$$\lim_{x \to +\infty}{y} = -\infty$$
\item Đồ thị:
	\begin{figure}[H]
		\centering
		\includegraphics[width=6cm]{images/cover/dhsp_logo.png}
		\caption{Hình vẽ đồ thị bài 2}
	\end{figure}
	
\end{itemize}

